% ============================================================================
% Theorem 10': The Verifiability Trichotomy (completeness of the SVNN condition)
% Corrects the overstrong "MLPs admit no design-time bounds" claim.
% ============================================================================

\subsection{The Verifiability Trichotomy: Why SVNN Is Conjecturally the Unique Optimum}
\label{sec:trichotomy}

Proposition~\ref{prop:mlp_negative} established that standard MLPs do not
admit the KAN-style de~Boor curvature bound. That claim is correct but
\emph{incomplete}: it does not imply that MLPs admit \emph{no} closed-form
design-time bounds. ReLU-MLPs admit first-order bounds ($\varepsilon \leq
M_1 h/2$ with $M_1 = 1$ for ReLU---Lipschitz LUT error, computable and
sound), and tanh/sigmoid-MLPs are $C^\infty$ with global closed-form $M_2$,
hence \emph{satisfy} the SVNN curvature condition. What coupled
architectures genuinely lack is the \emph{joint} attainment of second-order
rate, linear LUT cost, closed-form sharp constants, and freedom from
reachability computation. The correct boundary is a \emph{trichotomy}, and
the SVNN condition is the regime where all four properties coexist
(sufficiency proven; the necessity direction is conjectural---item~5 below).

\setcounter{theorem}{9}
\begin{theorem}[The Verifiability Trichotomy]
\label{thm:trichotomy}
Consider LUT-compiled feedforward networks with gates from a gate set
$\mathcal{G}$, bounded weights, compact input domain. Define three
properties:
\begin{itemize}
    \item[(P1)] \textbf{Closed-form design-time bound}, computable in
    polynomial time in $(L, d, N)$ with $O(L d^2)$ per-evaluation cost;
    \item[(P2)] \textbf{Second-order rate}: LUT error $\leq c_k M_k h^k$
    with $k \geq 2$ and sharp constants (Theorem~\ref{thm:sharp-lut});
    \item[(P3)] \textbf{Tightness without reachability}: constants depend
    only on global gate data and weight norms---no reachable-set
    computation.
\end{itemize}

\begin{enumerate}
    \item \textbf{Regime 1 (SVNN): all three hold.} If every gate is
    univariate, element-wise, and $C^2$ on the validated domain with
    design-time-computable $M_k$ (Condition~2 of
    Theorem~\ref{thm:svnn_conditions}, stated on the validated input
    domain $X$), and the Box-Continuation condition holds
    (Lemma~\ref{lem:box_continuation}: per-layer reachable boxes stay
    within the LUT table domain), then (P1)+(P2)+(P3) hold with the
    sharp constants of Theorem~\ref{thm:sharp-lut}. This is the SVNN
    condition with the domain-continuation refinement: gates with
    domain-dependent curvature (e.g., $e^x$) satisfy Conditions~1--2
    on a validated box but \emph{not} P3, since $M_k$ then depends on
    the reachable set---they belong to Regime~4 unless per-layer
    domains are certified.

    \item \textbf{Regime 2 (Lipschitz-only gates, e.g.\ ReLU): P1+P3 hold,
    P2 fails.} The minimax $N$-point approximation error of the $C^1$
    (Lipschitz) class is $\Theta(1/N)$ (classical width result); the
    kink cell forces error $\geq c \cdot M_1 \cdot h$, so the rate is at
    most $h$---first-order. ReLU-MLPs therefore \emph{do} admit closed-form
    first-order bounds (correcting the overstrong dichotomy), but not
    second-order ones. Measured: slope $-1.034$ (E66), vs $-2.04$ for
    $C^2$ activations.

    \item \textbf{Regime 3 (multivariate gates, e.g.\ products,
    softmax, normalization): linear LUT cost fails; tightness NP-hard.}
    Tensor-product LUTs cost $N^m$ per gate---exponential in the arity
    $m$ for unbounded-arity gates such as softmax or layer normalization
    (measured: storage $N^2$ for binary gates, $N^4$ for quaternary,
    E66). The per-dimension second-order rate \emph{survives}
    (measured slopes $-2.11$, $-1.82$, E66), so the failure is not in
    P2's rate but in P1's linear-cost requirement and in \emph{tight
    propagation}: computing tight constants requires solving
    polynomial-optimization reachability---NP-hard even for arity-2
    product gates (the 3-SAT embedding of Theorem~\ref{thm:necessity},
    Eq.~\ref{eq:product_gate}).

    \item \textbf{Regime 4 (domain-dependent curvature, e.g.\ $e^x$,
    $x \sin x$): P3 fails.} $M_k$ is finite on the reachable set but
    depends on it; closed-form tight certification would require the
    reachable set itself---reachability, NP-hard in general. The
    polynomial-time alternative (interval propagation on a global
    superset box) is sound but exponentially loose in $L$ (classic
    interval explosion: constants $\|W\|_\infty^L$). The Box-Continuation
    Lemma provides the certified escape: per-layer box propagation
    \emph{is} polynomial for SVNN-structured nets and restores P3.

    \item \textbf{Necessity (completeness, conditional).} Under
    $\mathrm{P} \neq \mathrm{NP}$ ($\exists\mathbb{R} \neq \mathrm{P}$),
    any family satisfying (P1)+(P2)+(P3) must lie in Regime~1: (P2)
    excludes first-order-rate gates---for kinked (ReLU-type) activations
    the kink cell forces interpolation error $\geq c\,M_1 h$ regardless
    of $N$, so such classes are not rate-2 recoverable by any
    design-time-computable affine (piecewise-linear) scheme, and
    Regime~2 is excluded. Note that the precise rate-2 boundary is
    bounded second variation (kink-free $C^1$, not necessarily $C^2$):
    the honest claim is the exclusion of classes that are \emph{not}
    rate-2 recoverable, with Condition~2 of
    Theorem~\ref{thm:svnn_conditions} certifying the sufficient $C^2$
    form with design-time-computable $M_k$. (P3) forces
    design-time-computable $M_k$ independent of the reachable set
    (else Regime~4) and univariate gates (else NP-hard tightness,
    Regime~3); (P1)-with-linear-cost forces univariate gates (else
    $N^m$). We state this necessity direction as a \emph{conjecture
    with strong evidence}---the three regime exclusions are each
    supported by the corresponding hardness/rate results above, but a
    fully formal necessity proof would require making the
    ``design-time-computable'' notion a formal complexity class, which
    we leave as an open problem.
\end{enumerate}
\end{theorem}

\vspace{2pt}
\noindent\textit{Proof sketch.}
\textbf{(1)} Theorem~\ref{thm:sharp-lut} + linear propagation.
\textbf{(2)} Lower bound on the $C^1$ width: for any $N$-point piecewise
scheme, a tent function of slope $M_1$ forces error
$\geq c M_1 (b-a)/N$ (the kink-cell argument suffices). Empirically
(E66): slope $-1.034$ for ReLU vs $-2.04$ for tanh on $[-3,3]$, $R^2 >
0.999$. \textbf{(3)} The product-gate hardness is the self-contained
reduction of Theorem~\ref{thm:necessity}; unbounded-arity gates make
even the LUT itself exponential ($N^m$). \textbf{(4)} For $e^x$-type
gates, $M_k = e^b$ where $b = \sup$ of the reachable set---computing $b$
is the reachability problem; the closed-form surrogate (global box from
weight norms) has constants exploding in $L$. \textbf{(5)} Combine:
(P2) $\Rightarrow$ rate-2 recoverability (Regime 2 excluded: the kink
cell of ReLU-type gates forces first-order error $c M_1 h$; the
certified sufficient form is $C^2$ with design-time-computable $M_k$,
though bounded second variation already yields the rate), (P3)
$\Rightarrow$ global computable $M_k$ and univariate (Regimes 3--4
excluded), (P1) $\Rightarrow$ linear-cost LUT (Regime 3's $N^m$
excluded). Hence Regime 1. $\square$

\vspace{2pt}
\noindent\textbf{Consequences.}
\textit{(i) Honest correction of the SVNN/MLP dichotomy.} The paper's
earlier claim that ``standard MLPs do not admit design-time bounds''
(Proposition~\ref{prop:mlp_negative}) is refined: ReLU-MLPs admit
\emph{first-order} bounds, tanh-MLPs \emph{satisfy} the SVNN curvature
condition; what MLPs lack is the joint attainment of all four properties.
\textit{(ii) Design guidance.} The trichotomy tells the compiler designer
exactly which gate sets are certifiable at second order, which at first
order, and which not at all---converting the SVNN condition from a
sufficient criterion into the \emph{characterization} of the feasible
design space. \textit{(iii) The unique-optimum reading.} SVNN is not
``one verifiable family among many'': it is \emph{conjecturally} the
unique regime where second-order rate, linear LUT cost, closed-form
sharp constants, and reachability-freedom coexist---the necessity
direction of the trichotomy is a conjecture with strong evidence, and
turning it into a theorem (formalizing ``design-time-computable'') is
an open problem. \textit{(iv) Empirical confirmation (E66).}
Rate trichotomy measured: $\{-2.04, -2.06, -1.77\}$ for tanh/sin/
B-spline (Regime 1), $-1.03$ for ReLU (Regime 2); tensor-LUT storage
$N^m$ for arity $m$ (Regime 3); tanh constant ratio $0.974 \times
M_2 h^2/8$ (sharpness of the Regime-1 constant).
